% ============================================================
% 分片：s05a2-comment  对应原文 PDF 第 58 页（印刷页 58）前 21 行（5.3 标题之前）
% 负责章节：§5.2 Hypersurfaces In Toric Varieties 的收尾
%           （Batyrev 句续文 + 无编号小标题 "A General Comment" 及其两段）
% 译者：t07-s34   审校：t03-s2b（逐式核对完成）   审校状态：二审
% ============================================================
% 边界说明（Lead 裁定）：
%   • 本文件首句紧接 sections/s05a-wps-toric.tex 末句“……值得注意的是，Batyrev [40] 给出了一类”，
%     是其续文；两文件之间不要插入空行或分段，以保证句子无缝拼接。
%   • 本文件止于原文 5.3 标题（HYPERSURFACES IN GRASSMANNIANS）之前；
%     5.3 及其后正文由 sections/s05b-grassmann-intersect.tex 接手。
%   • 原文 "A General Comment" 为无编号下划线小标题，按 SPEC §3 记作 \subsubsection*{}。
%   • 原文粗体 $\mathbb{C},\mathbb{R}$ 按 GLOSSARY/SPEC 4.5 统一为 $\mathbb{C},\mathbb{R}$；
%     辛商记号原排作“//”，本文件写作 $/\!/$。
% ============================================================
% ==== 原文 PDF 第 58 页（印刷页 58）====
\label{op:58}
% [跨页接续: 原文第 57 页末至第 58 页首；本文件首句为 s05a 末句的续文，构成完整句：
%  "It is worth noting that Batyrev [40] gave a kind / of Landau-Ginzburg description of the
%   cohomology of a general toric hypersurface; the facts we have just explained presumably give
%   a field theoretic setting for his construction."]
Landau--Ginzburg 描述，以刻画一般环面超曲面的上同调；我们刚刚说明的这些事实，想必为他的构造提供了一个场论背景。

\subsection*{一点一般性评论}

一个 $N-d$ 维的光滑环面簇 $Z$，若 $H^2(Z,\mathbb{R})$ 是 $d$ 维的，就总可以作为 $\mathbb{C}^N$ 对 $U(1)^d$ 的辛商得到。辛商 $\mathbb{C}^N/\!/U(1)^d$ 依赖于 $d$ 个参数 $r_1,\dots,r_d$，即可以加到 $D$ 函数上的那些常数项。按照 Duistermaat--Heckman 定理，从一个固定的 $\mathbb{C}^N$ 辛结构出发，$Z=\mathbb{C}^N/\!/U(1)^d$ 上自然诱导的辛结构随 $r_i$ 线性地变化。换言之，在 $H^2(Z,\mathbb{R})$ 的某个基 $e_i$ 下，$Z=\mathbb{C}^N/\!/U(1)^d$ 的凯勒类为 $[\omega]=\sum_i r_ie_i$。
% [Tutor Note: Duistermaat--Heckman 定理在此用到的确切内容是：对带 Hamiltonian $G$-作用的辛流形，约化空间 $\mu^{-1}(r)/G$ 上约化辛形式的上同调类随矩映射水平 $r$（Lie 代数对偶中的常数）线性变化。$\mathbb{C}^N$ 上的线性 $U(1)^d$ 作用其矩映射是二次型（即 $D$ 函数），把常数项 $r_i$ 平移过去只改变矩映射水平，故 $Z$ 的辛形式（即凯勒形式）的类可写成 $[\omega]=\sum_i r_ie_i$。这正是本文把 Fayet--Iliopoulos 参数 $r_i$ 当作 $Z$ 的凯勒类坐标的数学依据，也是 §5.2 前面各例中“$r_j$ 决定凯勒类”的一般化。]

现在考虑 Calabi--Yau 超曲面 $X\subset Z$。$H^2(X,\mathbb{R})$ 有一些生成元可以通过把 $Z$ 的上同调限制到 $X$ 上而得到。然而 $H^2(Z,\mathbb{R})$ 的限制可能只给出 $H^2(X,\mathbb{R})$ 的一部分（正如 [7] 中所研究的例子）。在这种情形下，$X$ 上最一般的凯勒度规并不便于通过嵌入 $Z$ 来获得，而 $r_i$ 只是参数化 $X$ 的凯勒类的自然线性坐标的一个子集。当这种情况发生时，我们所构造的相图并不是靶空间为 $X$ 的 $\sigma$ 模型的完整相图，而是它在 $r_i$ 张成的子空间上的限制。
% [Tutor Note: 限制映射 $i^*:H^2(Z,\mathbb{R})\to H^2(X,\mathbb{R})$（$i:X\hookrightarrow Z$ 为嵌入）未必满射。注意一个容易搞错的点：若 $X\subset Z$ 是**丰沛（ample）**超曲面且 $\dim_{\mathbb C}X\ge2$，Lefschetz 超平面定理说 $i^*$ 在低阶上其实是同构（$H^2$ 亦然），此时 $H^2(X)$ 的每个类都来自环境空间。原文“只给出一部分”的情形（[7] 的例子）之所以出现，是因为超曲面类并非总能取成 ample、或 $Z$ 自身带有非平凡 $H^2$ 之外的几何（例如 $X$ 的凯勒锥含有 $i^*$ 像之外的方向）；此时 $r_i$ 只遍历 $H^2(X,\mathbb R)$ 的一个子空间。物理上这意味着线性 $\sigma$ 模型（只有 $d$ 个 FI 参数）无法覆盖 $X$ 的全部凯勒模，故本文构造的相图只是完整相图在该子空间上的切片——这一点在 §5.5 讨论拓扑改变（凯勒锥的不同锥面）时必须牢记。]

% 本文件止于 5.3 标题之前；其后的 5.3/5.4 由 sections/s05b-grassmann-intersect.tex 接手。
